\tzpanchor{0239}
\tzpentryheader{0239}{Buckingham $\Pi$-Theorem Application}

\begingroup\small
\begingroup\scriptsize
\setlength{\tabcolsep}{4pt}
\renewcommand{\arraystretch}{0.92}
\noindent\begin{tabularx}{\linewidth}{@{}p{0.24\linewidth}p{0.36\linewidth}X@{}}
\textbf{UUID} & \multicolumn{2}{l@{}}{\tzpcode{5c93317d-80ed-5814-9f32-c5c9d7a6ca66}}\\
\textbf{PiID} & \tzpcode{5271201909145} & \textbf{TZPi}\quad \tzpcode{1}\quad \textbf{PiPE}\quad \tzpcode{246}\\
\textbf{RTEID} & \textbf{Poloidal $\theta$}\quad \tzpcode{928.0901 rad} & \textbf{Toroidal $\phi$}\quad \tzpcode{01.5292 rad}\\
\textbf{TZPID} & \tzpcode{ID0239} & \textbf{Node Dependencies}\quad \tzpidref{0238}\\
\textbf{Registry} & \tzpcode{registry\_id\_0239} & \textbf{Scale}\quad transfer\\
\textbf{LOC Classification} & QA641 manifolds and topology & QC174.17 mathematical physics\\
\textbf{arXiv Classification} & Primary: math-ph & Cross-list: gr-qc, math.DG\\
\textbf{Primary Support} & Variable Count & Dimension Count\\
\textbf{Closure Support} & Pi Groups & \\
\end{tabularx}\par
\endgroup
\endgroup

\tzpsection{Canonical Statement}
Dimensional analysis applied to gyromagnetic systems with fundamental quantities mass M, length R, time period T, magnetic field B, magnetic moment mu, electrical conductivity sigma and viscosity eta yields dimensionless groups including the magnetic Reynolds number R\_m and Elsasser number Lambda.

\tzpsection{Canonical Equation}
\[
R_m=\mu_0\sigma vL
\]
\[
R_m=\mu_0\sigma L^2\omega
\]

\begin{minipage}[t]{0.49\linewidth}
\tzpsection{Canonical Symbol Key}
\begingroup
\small
\raggedright
\textbf{Source equation symbols.}\par
\begin{itemize}[leftmargin=1.35em,itemsep=1pt,topsep=1pt,parsep=0pt]
    \item $R_m$: source equation symbol.
    \item $\mu_0$: vacuum permeability constant.
    \item $\sigma$: conductivity or coupling coefficient in the source equation.
    \item $L$: characteristic length scale or Lagrangian carrier in the entry relation.
    \item $\omega$: angular frequency term in the source equation.
\end{itemize}
\endgroup
\end{minipage}\hfill
\begin{minipage}[t]{0.47\linewidth}
\tzpsection{Wolfram Form}
\begingroup
\scriptsize
\raggedright
\textbf{Source Wolfram model.}\par
\begin{Verbatim}[fontsize=\tiny,breaklines=true,breakanywhere=true]
(* TZPID ID0239 generated source Wolfram model *)
ClearAll[K0239, Cr0239, T, R, A, W, I, N];
eq0239$1 = HoldForm[RM == mu0*sigma*vL];
eq0239$2 = HoldForm[RM == mu0*sigma*L^2*omega];
eq0239$3 = HoldForm[RM == vL/eta];

K0239 = HoldForm[{eq0239$1, eq0239$2, eq0239$3}];

N = "normalized TRAWIN closure object for Buckingham $\Pi$-Theorem Application";
I = "Pi Groups integration/comparison layer";
W = "Dimension Count wave or operator carrier";
A = "Variable Count admissible action/amplitude condition";
R = "Dimension Count rotation/curl preservation layer";
T = "Variable Count local source condition";

Cr0239[expr_] := HoldForm[N[I[W[A[R[T[expr]]]]]]];

Result0239 = Cr0239[K0239];
\end{Verbatim}
\endgroup
\end{minipage}

\par\vspace{0.45\baselineskip}
\noindent\begin{minipage}[t]{\linewidth}
\begingroup
\small
\raggedright
\textbf{TRAWIN Operator Key.}\par
\noindent\textbf{Time/localization operator:} $\mathsf{T}\!\left[\mathrm{local}\right]$ Variable Count local source condition.\par
\noindent\textbf{Rotation/curl operator:} $\mathsf{R}\!\left[\nabla\times\right]$ Dimension Count rotation/curl preservation layer.\par
\noindent\textbf{Action/amplitude operator:} $\mathsf{A}\!\left[\mathrm{admissible}\right]$ Variable Count admissible action/amplitude condition.\par
\noindent\textbf{Wave/d'Alembertian operator:} $\mathsf{W}\!\left[\Box\right]$ Dimension Count wave or operator carrier.\par
\noindent\textbf{Integration operator:} $\mathsf{I}\!\left[\int\right]$ Pi Groups integration/comparison layer.\par
\noindent\textbf{Normalization operator:} $\mathsf{N}\!\left[\mathrm{normalized}\right]$ normalized TRAWIN closure object for Buckingham \$\textbackslash{}Pi\$-Theorem Application.\par
\endgroup
\end{minipage}

\clearpage
\noindent\textbf{[ID 0239] Buckingham $\Pi$-Theorem Application}\par
\vspace{0.35em}
\tzpsection{Derivation Layer}
\paragraph{Primary Equation.}
\begin{equation}
\frac{\mathbf{B}^2}{2\mu}
=
\frac{1}{2}\rho(\Omega r)^2 .
\end{equation}

\paragraph{Primary Derivation.}
The magnetic energy density is:

\begin{equation}
u_B
=
\frac{\mathbf{B}^2}{2\mu}.
\end{equation}

The rotational kinetic energy density associated with angular velocity $\Omega$ at radius $r$ is:

\begin{equation}
u_{\Omega}
=
\frac{1}{2}\rho v^2 .
\end{equation}

Since tangential velocity is:

\begin{equation}
v=\Omega r,
\end{equation}

we obtain:

\begin{equation}
u_{\Omega}
=
\frac{1}{2}\rho(\Omega r)^2 .
\end{equation}

Equipartition requires:

\begin{equation}
\boxed{
\frac{\mathbf{B}^2}{2\mu}
=
\frac{1}{2}\rho(\Omega r)^2 .
}
\end{equation}

\paragraph{Interpretation.}
Magnetic--Coriolis equipartition identifies the balance point where magnetic field energy and rotational kinetic energy density match.

\par\vspace{0.35em}
\noindent\textbf{Derivable Consequences.}\quad
\begingroup
\footnotesize
\raggedright
The canonical equation anchors Buckingham \$\textbackslash{}Pi\$-Theorem Application by connecting the source condition to the derivation layer and proof certificate.
\par\endgroup

\tzpsection{Equation Support}
\noindent\textbf{1. Variable Count}\par
\[
R_m=\mu_0\sigma vL
\]
\noindent This fixes the number of governing variables in the dimensional analysis setup.\par
\noindent\textbf{2. Dimension Count}\par
\[
R_m=\mu_0\sigma L^2\omega
\]
\noindent This fixes the rank of the independent dimensional basis.\par
\noindent\textbf{3. Pi Groups}\par
\[
R_m=\frac{vL}{\eta}
\]
\noindent This closes the entry with the resulting number of independent Buckingham Pi groups.\par

\clearpage
\noindent\textbf{[ID 0239] Buckingham $\Pi$-Theorem Application}\par
\vspace{0.35em}
\tzpsection{Dictionary}
\begingroup\small
\tzpsection{Definition}
Buckingham Pi-Theorem Application is a registry definition in Information Synchronization with secondary pressure from Field Theory and Action Principles.

% BEGIN TZP CONTEXTUAL MEANING
\tzpsection{Contextual Meaning}
\begingroup\footnotesize\setlength{\emergencystretch}{3em}\sloppy
\textbf{Formal statement:} Dimensional analysis applied to gyromagnetic systems with fundamental quantities mass M, length R, time period T, magnetic field B, magnetic moment mu, electrical conductivity sigma and viscosity eta yields dimensionless groups including the magnetic Reynolds number R sub m and Elsasser number Lambda. \textbf{Contextual meaning:} The entry reads Buckingham Pi-Theorem Application as a controlled technical headword whose equation layer, proof route, and registry role are interpreted together. \textbf{Derivable consequences:} The entry now states the Buckingham count cleanly instead of mixing prose and inline dollar math inside display mode. \textbf{Lemma support:} Variable Count; Dimension Count; Pi Groups.\par
\endgroup
% END TZP CONTEXTUAL MEANING

\tzpsection{Technical Interpretation}
Magnetic--Coriolis equipartition identifies the balance point where magnetic field energy and rotational kinetic energy density match.

\tzpsection{Equation Grounding}
Equation anchors: R sub m=mu sub 0sigma vL; R sub m=mu sub 0sigma L to the power 2omega; and R sub m=ratio vL eta. Proof route: show that seven governing quantities with four independent dimensions yield exactly three dimensionless Pi groups. Formalization route: prove the Pi-group count from the stated dimensional-rank relation.

\tzpsection{Interpretive Role}
The entry now states the Buckingham count cleanly instead of mixing prose and inline dollar math inside display mode.
\endgroup

\tzpsection{TZP Thesaurus}
\begin{enumerate}[leftmargin=1.6em,itemsep=6pt,topsep=4pt]
\item \textbf{Dimensional Matrix Reduction Framework}: The rigorous algebraic method of reducing a set of complex physical variables into a smaller set of fundamental, dimensionless ratios.
\item \textbf{Dimensionless Parameter Group Extraction}: The systematic generation of the core numbers (like Reynolds or Elsasser numbers) that truly govern the system's scale-invariant physics.
\item \textbf{Scaling Law Algebraic Synthesis}: The formulation of universal physical equations strictly utilizing combinations of unitless Pi-groups.
\end{enumerate}

\tzpsection{Encyclopedia Mathematica}
\begingroup\footnotesize\sloppy
The visual plate below is reserved for the source-specific equation and progression graphic for [ID 0239]. It is included only as a companion to the canonical equation, not as a substitute for the formal text.
\par\endgroup
\ifTZPDarkEdition
\tzpdictionaryvisualband{D:/TZPID/kdp_workflow/build/tikz_figures/ID0239_dictionary_figure_dark.pdf}
\fi

\clearpage
\begingroup\tzpsupportsetup
\noindent\textbf{\fontsize{10.8pt}{12.8pt}\selectfont [ID 0239] Constructive Logic and Proof Certificate}\par
\noindent\textbf{\fontsize{10pt}{12pt}\selectfont Buckingham $\Pi$-Theorem Application}\par
\vspace{0.45em}
\begingroup
\footnotesize
\setlength{\parskip}{0.22em}
\setlength{\parindent}{0pt}
\noindent\textbf{Curry--Howard--Lambek Interpretation}\par
Under the Curry--Howard--Lambek correspondence, this registry entry is read as a typed proof
object: the stated source condition supplies the proposition, the derivation supplies the
morphism, and the proof certificate records the machine handles needed to check the obligation
in the formal registry.

\vspace{0.45em}
\noindent\textbf{CHL Proof}\par
\textbf{Statement:} dimensional analysis applied to gyromagnetic systems with fundamental quantities mass M,
length R, time period T, magnetic field B, magnetic moment...

\textbf{Logic Validation:}\\
\textbf{Proposition:} ``Buckingham \textbackslash{}Pi-Theorem Application is valid as a formal registry statement when its
source equation and semantic claim are read together.''\\
\textbf{Proof:} The source semantics state that dimensional analysis applied to gyromagnetic systems
with fundamental quantities mass M, length R, time period T, magnetic field B, magnetic
moment mu, electrical conductivity sigma and viscosity eta yields dimensionless groups
including the magnetic Reynolds number R\textbackslash{}\_m and Elsasser number Lambda. Magnetic--
Coriolis equipartition identifies the balance point where magnetic field energy and
rotational kinetic energy density match. The entry now states the Buckingham count
cleanly instead of mixing prose and inline dollar math inside display mode.. The CHL
validation treats the equation as the formal anchor and the semantic text as the
interpretation constraint.

\textbf{VERDICT:} \ensuremath{\checkmark}\ \textbf{TRUE} --- source-backed formal validation

\textbf{Formal Status:} \ensuremath{\checkmark}\ THEORETICAL -- source-backed CHL proof obligation

\textbf{Physical Status:} THEORETICAL -- requires model-specific empirical interpretation
\par\vspace{0.45em}
\noindent{\tiny\textbf{Isabelle/HOL.} \tzpplaincode{physics\_validity\_from\_model, external\_validity\_package, registry\_number\_matches, equation\_count\_matches}}\par
\ifTZPDarkEdition
\tzpproofvisualband{D:/TZPID/kdp_workflow/build/tikz_figures/ID0239_proof_certificate_figure_dark.pdf}
\fi
\endgroup
\endgroup
